\subsection{GloVe: Global Vectors for Word Representation}\label{sec:glove-global-vectors-for-word-representation}

\subsubsection{Motivation behind GloVe}\label{sec:motivation-behind-glove}

After the success of \textbf{Word2Vec}, researchers observed that its embeddings exhibited remarkable semantic properties:
\begin{itemize}
    \item Semantically related words were close together;
    \item Vector arithmetic captured analogies such as:
    \begin{equation*}
        \emph{king} - \emph{man} + \emph{woman} \approx \emph{queen}
    \end{equation*}
    \item Relationships such as \emph{country}-\emph{capital} emerged naturally.
\end{itemize}
However, \textbf{Word2Vec never explicitly models these relationships}. Instead, they appear as a by-product of the prediction task. During training, the model learns to predict surrounding words from a target word (Skip-Gram) or vice versa (CBOW), and the \hl{semantic structure of the embedding space emerges implicitly.}

\highspace
\textcolor{Green3}{\faIcon{question-circle} \textbf{The Central Question.}} The authors of \definition{GloVe (Global Vectors for Word Representation)} \cite{pennington2014glove} wondered whether \hl{these semantic regularities could be learned more directly.} Instead of learning embeddings by solving a prediction problem, they proposed \hl{learning them from the statistical structure of the entire corpus.} In other words:
\begin{itemize}
    \item Word2Vec asks: \emph{Can we predict neighboring words?}
    \item GloVe asks: \emph{\textbf{What do the co-occurrence statistics of all words tell us about their meanings?}}
\end{itemize}
This is the main conceptual difference between the two approaches.

\highspace
\textcolor{DarkOrange3}{\faIcon{balance-scale} \textbf{Local versus global information.}} Recall how Word2Vec learns. Suppose we observe the sentence:
\begin{center}
    \emph{The cat sits on the mat.}
\end{center}
When training on the word \emph{cat}, Word2Vec only looks at the words inside a small context window, for example:
\begin{center}
    \emph{The} \quad \emph{cat} \quad \emph{sits}
\end{center}
Later it processes another sentence:
\begin{center}
    \emph{The cat drinks milk.}
\end{center}
Using another local window. Therefore, \textbf{each update depends only on a small portion of the corpus}. Although millions of these local updates eventually produce excellent embeddings, the \hl{model never explicitly computes how often every pair of words co-occurs in the entire corpus.}

\newpage

\begin{flushleft}
    \textcolor{Green3}{\faIcon[regular]{lightbulb} \textbf{GloVe's idea}}
\end{flushleft}
GloVe starts from a completely different observation. Instead of asking ``\emph{which words appear around ice?}'', it asks \emph{\textbf{how often does ice appear together with every other word in the corpus?}}. For example, suppose we count co-occurrences in a large corpus and find:
\begin{center}
    \begin{tabular}{@{} l l @{}}
        \toprule
        \textbf{Context word} & \textbf{Times observed with \emph{ice}} \\
        \midrule
        \emph{solid} & 1900 \\
        \emph{water} & 3000 \\
        \emph{cold}  & 2500 \\
        \emph{steam} & 22 \\
        \emph{fashion} & 17 \\
        \bottomrule
    \end{tabular}
\end{center}
\noindent
These counts immediately reveal something important:
\begin{itemize}
    \item \emph{ice} is strongly related to \emph{solid};
    \item \emph{ice} is moderately related to \emph{water};
    \item \emph{ice} is almost unrelated to \emph{fashion}.
\end{itemize}
Instead of predicting neighboring words one example at a time, \textbf{GloVe attempts to directly learn embeddings that explain these global statistics}.

\highspace
\textcolor{Green3}{\faIcon{question-circle} \textbf{Why is this useful?}} The intuition is that \hl{meaning is reflected by patterns of co-occurrence.} For example:
\begin{itemize}
    \item \emph{ice} frequently appears with \emph{cold}, \emph{snow}, \emph{winter}, and \emph{frozen};
    \item \emph{steam} frequently appears with \emph{hot}, \emph{boiling}, \emph{vapor}, and \emph{heat}.
\end{itemize}
Notice something interesting: both words occur with \emph{water}, but they differ greatly in how often they occur with \emph{solid} or \emph{gas}. Therefore, simply knowing that two words co-occur is not sufficient. The \textbf{relative frequencies} of their co-occurrences contain richer semantic information. This observation is the key insight that motivates GloVe.

\highspace
The GloVe authors argued that semantic relationships should not simply emerge by chance during prediction. Instead, they should be encoded directly from the statistical properties of the corpus. As stated in the original paper:
\begin{itemize}
    \item \textbf{Word2Vec implicitly} learns semantic vector offsets through prediction;
    \item \textbf{GloVe explicitly} models the statistical relationships that give rise to those offsets.
\end{itemize}

\highspace
In summary, the motivation behind GloVe is to combine the strengths of:
\begin{itemize}
    \item \textbf{Count-based methods}, which exploit the \textbf{global co-occurrence\break statistics} of the corpus;
    \item \textbf{Prediction-based methods} like Word2Vec, which produce dense, high-quality embeddings.
\end{itemize}
Rather than learning word vectors only through local prediction tasks, GloVe directly learns embeddings that capture the \textbf{global statistical structure of language}. This idea will naturally lead to the next section, where we introduce the global co-occurrence matrix and explain how GloVe uses it to learn word embeddings.

\highspace
The \textbf{motivation} behind \textbf{GloVe (Global Vectors for Word Representation)} is to learn word embeddings by exploiting the \textbf{global statistical structure} of the corpus. While \textbf{Word2Vec} learns embeddings \emph{implicitly} by solving a local prediction task (predicting surrounding words), GloVe learns them \emph{explicitly} by modeling the global co-occurrence statistics of words. Rather than asking \emph{``Which word should be predicted next?''}, GloVe asks \emph{``Which word pairs frequently occur together throughout the entire corpus?''}.